% frontmatter/how_to_use_neuroscience_analogies.tex
\chapter*{如何阅读本书的神经科学类比？}
\addcontentsline{toc}{chapter}{如何阅读本书的神经科学类比？}

在阅读本书时，你会发现我们频繁且深入地使用了脑科学与神经科学的现象来作为理解人工智能系统的直觉辅助。
在进入具体的章节之前，本指南将明确告诉你在本书中，\textbf{什么能用类比来理解，什么必须用数学来判断}。

\section*{为什么要用神经科学类比？}
现今的复杂人工智能推断、基于注意力的提取或是隐状态梯度调整，对于初学者常常是一堆孤立的矩阵公式。
但是，在自然界中，大脑作为一个极其复杂、能量受限、拥有长短期记忆分布和突触级局布可塑性的动态系统，
正是我们在工程中试图去模拟其“信息处理与时序行为表现”的对象。

我们使用由于神经科学积累的名词（例如工作记忆、海马体长程固化、突触可塑性等），是因为：
\begin{enumerate}
    \item 它们提供了直截了当的\textbf{计算角色}和\textbf{时序特征}预设。
    \item 许多顶尖人工神经网络架构的源头灵感常常来源于粗燥的仿生直觉。
\end{enumerate}

\section*{类比的三层结构阅读法}
本书在触及任何一个核心模型时，我们严格采用了三层解析法。阅读时请始终保持以下心法：
\begin{itemize}
    \item \textbf{第一层：数学与算法形式化。}这是唯一真实发生的计算！无论是注意力还是 TTT，它们是 $A \times B + C$ 的浮点数截断行为。请用算法工程师的严苛去抠这些项。
    \item \textbf{第二层：架构与系统的动机。}理解为什么这些数学项组合在一起能够实现数据的关联筛选或测试时的二次回归。
    \item \textbf{第三层：特定生物学类比（辅以严格的边界排谬）。}这是在你脑海中为其附着宏观概念的标签。
\end{itemize}

\section*{类比必定会失效的几个红线}
在使用本书每一个 \textbf{【神经科学直觉】} 盒子时，请务必连同下方的 \textbf{【类比边界与失效条件】} 一起阅读。
你\textbf{绝不能}从本书的类比中推出以下错误结论：
\begin{enumerate}
    \item \textbf{“模型等价于这块神经组织”}：无论 TTT-MLP 或 Titans 上的动态优化写得多么像突触重塑，模型网络更新是由极其显式清晰、高度全局导向的反向传播计算目标控制的！这是强人类工程学控制。而真实生物突触更多依赖复杂的生化级联、局部电位和模糊缓慢的时序修正机制。
    \item \textbf{“我们已经懂大脑了” 或 “我们创造了数字大脑”}：我们所讲述的任何在线优化拟合方程，相比于大脑动辄数十种神经递质全天候处于异步耦合工作状态下的复杂拓扑网络，极大概率依然不过是管中窥豹的简化特征映射抽调罢了。
\end{enumerate}

本教材的核心理念是：\textbf{用严格的形式化保证真，用清楚的解释保证懂，用神经科学类比保证直觉，但永远不把类比偷换成结论。}
阅读全书时，请将类比视为脚手架（Scaffold），在完成数学逻辑大楼的理解后，你可以也应当随时拆除它。
